% Source pp. 20-26 (Lectures1-6.pdf).
\section{Casimir Effect}

To understand possible manifestations of the vacuum energy, let us now
briefly discuss the \emph{Casimir Effect}. Originally, the Casimir effect
refers to the situation where we have two parallel conducting plates, of large
size $L\times L$, separated by a distance $a$ (figure~\ref{fig:plates}).
\begin{figure}[ht]
\centering\figfile{plates}
\caption{Two parallel conducting plates of size $L\times L$ at separation $a$.}
\label{fig:plates}
\end{figure}
Due to the quantum nature of the electromagnetic field in between the plates,
more precisely, due to the effect of the plates on the zero-point energy of
quantum electromagnetic field, the plates are attracted with the force
\[
\frac{\text{Force}}{L^2}=-\hbar c\,\frac{\pi^2}{240\,a^4}\,,
\]
where it is assumed that
\[
L\gg a\gg\text{atomic distances.}
\]

\subsection*{Casimir effect in KG theory}

Let us study the nature of this effect using our scalar field theory, the KG
theory, instead of the electrodynamics; the result is not going to be very
different. Since the real photons are massless (and to make calculations
simpler), we will assume that $m^2=0$. In macroscopic electromagnetic theory
(valid at the scales $\gg$ atomic distances) the presence of the conducting
plates is taken into account by imposing the boundary conditions
$\mathbf{E}_\parallel=0$, $\mathbf{B}_\perp=0$ at the surface of the conductor.
To mimic this situation we consider instead the KG field in the presence of
two plates with the boundary condition
\[
\varphi(\xx)=0\qquad\text{at the surface.}
\]
Let us concentrate attention on the field between the plates. In order to
take into account the boundary conditions, in this case we must use, instead
of the functions $f_\pp(x)$ and $f^*_\pp(x)$ above, the following complete set
of solutions
\[
f_{\pp_\parallel,n}(\xx,t)=\sqrt{\frac2a}\;e^{i\omega_{\pp_\parallel,n}t-i\pp_\parallel\cdot\xx_\parallel}\,
\sin\left(\frac{\pi n}{a}\,x_\perp\right)\,,
\]
and the corresponding complex-conjugated functions $f^*_{\pp_\parallel,n}$.
Here $\xx_\parallel$ and $x_\perp$ are the components of $\xx$ parallel and
perpendicular to the plates, $\xx=(\xx_\parallel,x_\perp)$, and
\[
\omega_{\pp_\parallel,n}=\sqrt{\pp_\parallel^2+\frac{\pi^2n^2}{a^2}}
\]
with positive integer $n$,
\[
n=1,2,3,\dots\,.
\]
These functions solve the KG equation with $m^2=0$ as well as the boundary
conditions
\[
f_{\pp_\parallel,n}\big|_{x_\perp=0}=0\,,\qquad\text{and}\qquad
f_{\pp_\parallel,n}\big|_{x_\perp=a}=0\,.
\]
Exactly as we did in the infinite-volume case, we now introduce the operators
\[
\hat A_{\pp_\parallel,n}=\int_0^a \dd x_\perp\int \dd^2\xx_\parallel\,
\bigl[\hat\pi\,f_{\pp_\parallel,n}-\hat\varphi\,\dot f_{\pp_\parallel,n}\bigr]\,,
\]
and similarly for $\hat A^\dagger_{\pp_\parallel,n}$, with
$f^*_{\pp_\parallel,n}$ replacing $f_{\pp_\parallel,n}$. It is not difficult
then to show that
\[
[\hat A_{\pp_\parallel,n},\hat A^\dagger_{\pp'_\parallel,n'}]
=2\omega_{\pp_\parallel,n}\,\delta_{nn'}\,(2\pi)^2\,\delta(\pp_\parallel-\pp'_\parallel)\,,
\]
and to compute the Hamiltonian of the part of the system residing between the
plates,
\[
\hat H=\sum_{n=1}^\infty\int\frac{\dd^2\pp_\parallel}{(2\pi)^2}\,\frac14\,
\Bigl(\hat A^\dagger_{\pp_\parallel,n}\hat A_{\pp_\parallel,n}+\hat A_{\pp_\parallel,n}\hat A^\dagger_{\pp_\parallel,n}\Bigr)
+a\,L^2\,(\varepsilon-\varepsilon_0)\,,
\]
where the last term comes from vacuum energy counterterm
\[
(\varepsilon-\varepsilon_0)\int_0^a \dd x_\perp\int \dd^2\xx_\parallel=a\,L^2\,(\varepsilon-\varepsilon_0)\,.
\]
We still have to demand that all the operators $\hat A_{\pp_\parallel,n}$
nullify the ground state, since otherwise $\hat H$ is unbounded from below.
Therefore the vacuum eigenvalue of this Hamiltonian is
\[
E_{\text{vac}}=\frac{L^2}{2}\sum_{n=1}^\infty\int\frac{\dd^2\pp_\parallel}{(2\pi)^2}\,
\sqrt{\pp_\parallel^2+\frac{\pi^2n^2}{a^2}}\;\Phi\left(\frac{\pp_\parallel^2+\pi^2n^2/a^2}{\Lambda^2}\right)
-a\,L^2\,\varepsilon_0+a\,L^2\,\varepsilon\,,
\]
where I have introduced again the cutoff factor $\Phi$; without it this
expression is as divergent as it was in the absence of the boundaries.

Now, the term $-a\,L^2\,\varepsilon_0$ is the divergent part of the vacuum
energy counterterm. It is
\[
-a\,L^2\,\varepsilon_0=-\frac{L^2}{2}\,a\int\frac{\dd p_\perp}{(2\pi)}\int\frac{\dd^2\pp_\parallel}{(2\pi)^2}\,
\sqrt{\pp_\parallel^2+p_\perp^2}\;\Phi\left(\frac{\pp_\parallel^2+p_\perp^2}{\Lambda^2}\right)\,,
\]
where I have separated the $\dd p_\perp$ and $\dd^2\pp_\parallel$ parts of the
integral over $\dd^3\pp=\dd p_\perp\dd^2\pp_\parallel$. After the change of
variables
\[
p_\perp=\pi\,\tau/a
\]
it takes the form
\[
-\frac{L^2}{2}\int_0^\infty \dd\tau\int\frac{\dd^2\pp_\parallel}{(2\pi)^2}\,
\sqrt{\pp_\parallel^2+\frac{\pi^2\tau^2}{a^2}}\;\Phi\left(\frac{\pp_\parallel^2+\pi^2\tau^2/a^2}{\Lambda^2}\right)\,.
\]
We see that the interesting part of the ground-state energy (rather, of the
energy per unit area of the plates) can be written as
\[
\frac{E}{L^2}-a\,\varepsilon=\int\frac{\dd^2\pp_\parallel}{(2\pi)^2}\,
\left[\sum_{n=1}^\infty F_{\pp_\parallel}(n)-\int_0^\infty F_{\pp_\parallel}(\tau)\dd\tau\right]\,,
\]
where
\[
F_{\pp_\parallel}(\tau)=\frac12\,\sqrt{\pp_\parallel^2+\pi^2\tau^2/a^2}\;\Phi\left(\frac{\pp_\parallel^2+\pi^2\tau^2/a^2}{\Lambda^2}\right)\,.
\]
The $\sum F(n)-\int F(\tau)$ expression above can be handled with the help of
the Abel-Plana summation formula (see the Appendix to this section)
\begin{equation}\label{eq:AbelPlana}
\sum_{n=1}^\infty F(n)-\int_0^\infty F(\tau)\dd\tau
=-\frac12\,F(0)+i\int_0^\infty\frac{F(it+0)-F(-it+0)}{e^{2\pi t}-1}\dd t\,,
\end{equation}
which is valid provided $F(\tau)$ is analytic in the right half-plane of
complex $\tau$. The first term
\[
-\frac12\,F_{\pp_\parallel}(0)\ \to\ -\frac14\int\frac{\dd^2\pp_\parallel}{(2\pi)^2}\,
\sqrt{\pp_\parallel^2}\;\Phi(\pp_\parallel^2/\Lambda^2)
\]
does not depend on the separation $a$; it is interpreted as the surface
energy. It diverges, again signaling ambiguity. When there is a surface, a
constant surface term can be added to the action. We ignore this term since,
being independent of $a$, it does not contribute to the force.

The interesting contribution is extracted from the second term, which involves
\[
i\int_0^\infty\frac{F_{\pp_\parallel}(it+0)-F_{\pp_\parallel}(-it+0)}{e^{2\pi t}-1}\dd t\,.
\]
The shift ``$+0$'' in the arguments in \eqref{eq:AbelPlana} is important
since the function $F_{\pp_\parallel}(\tau)$ has square-root branch points at
the imaginary axis. Note that if it was regular, the integral would vanish in
view of the obvious $\tau\to-\tau$ symmetry of $F_{\pp_\parallel}(\tau)$. In
fact, one can take advantage of this symmetry to write the integrand as
\[
i\int_0^\infty\frac{F(it+0)-F(it-0)}{e^{2\pi t}-1}\dd t
=-2\int_{|\pp_\parallel|a/\pi}^\infty\frac{\Im F(it+0)}{e^{2\pi t}-1}\dd t\,,
\]
where the last form reflects the fact that the function
$F_{\pp_\parallel}(\tau)$ has the branch cut from $i|\pp_\parallel|a/\pi$ to
$i\infty$ (and from $-i|\pp_\parallel|a/\pi$ to $-i\infty$). We have, at
$|t|>a|\pp_\parallel|/\pi$,
\[
\Im F_{\pp_\parallel}(it+0)=\frac12\,\sqrt{\frac{\pi^2t^2}{a^2}-\pp_\parallel^2}\;
\Phi\left(\frac{\pp_\parallel^2-\pi^2t^2/a^2}{\Lambda^2}\right):=\frac12\,G_{\pp_\parallel}(t)\,.
\]
Note that the factor $(e^{2\pi t}-1)^{-1}$ makes this integral fast
convergent at large $t$, even without the cutoff factor $\Phi$. With large
$\Lambda\gg a^{-1}$ the cutoff factor has no effect, and can be dropped,
\[
G_{\pp_\parallel}(t)\to\sqrt{\frac{\pi^2t^2}{a^2}-\pp_\parallel^2}\,.
\]
The interesting contribution to the vacuum energy then is
\[
-\int\frac{\dd t}{e^{2\pi t}-1}\int_{|\pp_\parallel|<\pi t/a}\frac{\dd^2\pp_\parallel}{(2\pi)^2}\,
\sqrt{\frac{\pi^2t^2}{a^2}-\pp_\parallel^2}\,,
\]
where I have interchanged the order in which the integrals are taken. The
integral over $\pp_\parallel$ is evaluated in closed form
\[
\int_0^{\pi^2t^2/a^2}\frac{\dd p^2}{4\pi}\,\sqrt{\frac{\pi^2t^2}{a^2}-p^2}=\frac1{6\pi}\left(\frac{\pi t}{a}\right)^3\,.
\]
The contribution to the energy is
\[
-\frac{\pi^2}{6\,a^3}\int_0^\infty\frac{t^3\dd t}{e^{2\pi t}-1}=-\frac{\pi^2}{6\,a^3}\,\frac{\Gamma(4)\,\zeta(4)}{(2\pi)^4}\,,
\]
where I have used
\[
\int_0^\infty\frac{x^{s-1}\dd x}{e^x-1}=\Gamma(s)\,\zeta(s)
\]
($\zeta(4)=\pi^4/90$). Thus we find
\[
\frac{E}{L^2}=a\,\varepsilon-\frac12\,F(0)-\frac{\pi^2}{1440\,a^3}\,,
\]
where $F(0)$ stands for the surface term $\int\frac{\dd^2\pp_\parallel}{(2\pi)^2}F_{\pp_\parallel}(0)$.

The contribution $a\,\varepsilon$ is not too interesting: there is the same
term in the vacuum energy density outside the space between the plates, and
hence this term does not lead to any force. The interesting term
$-\pi^2/1440\,a^3$ shows that the plates attract; the force (per unit plate
area) is
\[
\frac{\text{Force}}{L^2}=-\frac{\partial}{\partial a}\left(-\frac{\pi^2}{1440\,a^3}\right)=-\frac{\pi^2\,\hbar c}{480\,a^4}\,.
\]
Let me stress that this result is completely independent of the form of
auxiliary cutoff factor $\Phi$; this is exactly what we want from quantum
field theory. Another useful observation is that in order to obtain this
result we don't really need to send $\Lambda$ to Planck's energies or
something, it suffices to have $a^2\gg1/\Lambda^2$.

In the real case of electromagnetic field the calculations can be done in very
similar way. The resulting force is twice as large, due to the fact that
photon field has two degrees of freedom for each $\pp$, as the photon has two
polarization states.

Few words about the $a$-independent term
\[
-\frac12\,F(0)=-\frac1{16\pi}\int_0^\infty \dd\xi\,\sqrt\xi\;\Phi(\xi/\Lambda^2)\,.
\]
It does not contribute to the force, but still needs some interpretation. It
is divergent, and it is easy to design special counter-term to absorb this
divergence. But it is instructive to think about real case of electromagnetic
field and real conducting plates made of atoms, conducting electrons, and all
that. In this case the cutoff represents the limit of validity of the
macroscopic electrodynamics, and the cutoff energy $\Lambda$ in this
expression should be taken around atomic energy scale. Having sufficient
knowledge about the microscopic structure of the conductor we could in
principle calculate the cutoff factor and hence this contribution in terms of
the atomic parameters. The nice feature of the above result is in that it
predicts some attractive force between the plates which is completely
independent of all these microscopic details. This points to another area
where quantum field theory applies, when it is not that we can't know the
short-distance physics, rather we don't want to know. In this area the
quantum field theory describes universal, i.e.\ not too sensitive to
microscopic details, properties of matter.

\subsection*{Appendix to Sect.~3: Abel-Plana summation formula}

Let $F(\tau)$ be a function analytic in the right half-plane, and decaying at
$\Re\tau\to+\infty$ sufficiently fast, so that the integral
$\int_0^\infty F(\tau)\dd\tau$ converges. Start with the integral
\[
I=\int_{C_0}\frac{F(\tau)\dd\tau}{e^{2\pi i\tau}-1}\,,
\]
where the integration is over the contour $C_0$ which encloses all positive
integer points $\tau=1,2,3,\dots$, as shown in the left part of
figure~\ref{fig:abelplana}.
\begin{figure}[ht]
\centering\figfile{abelplana}
\caption{Left: the contour $C_0$ encircles the poles at $\tau=1,2,3,\dots$
counterclockwise and passes to the right of $\tau=0$. Right: the deformed
contour. The part $C_+$ above the real axis is turned onto the positive
imaginary axis, and $C_-$ onto the negative imaginary axis. The small
quarter-circles $C'_\pm$ go around $\tau=0$.}
\label{fig:abelplana}
\end{figure}
The integrand has poles at the integer points $\tau=n$, with the residues
$F(n)/2\pi i$. Therefore
\[
I=\sum_{n=1}^\infty F(n)\,.
\]
On the other hand, the integral can be split into two parts
\[
I=\int_{C_0}=\int_{C_+}+\int_{C_-}=I_++I_-\,,
\]
where $C_+$ is the part of the contour $C_0$ that lies above the real axis,
and $C_-$ is the part below the real axis. Note that the factor
\[
\frac1{e^{2\pi i\tau}-1}
\]
decays fast into the lower half-plane, hence the contour $C_-$ can be rotated
into the lower half-plane to go along the contour $\tilde C_-$,
\[
I_-=\int_{\tilde C_-}\frac{F(\tau)\dd\tau}{e^{2\pi i\tau}-1}
=-i\int_\epsilon^\infty\frac{F(-it+0)\dd t}{e^{2\pi t}-1}+\int_{C'_-}\frac{F(\tau)\dd\tau}{e^{2\pi i\tau}-1}\,,
\]
where for the straight part of $\tilde C_-$ I have changed
\[
\tau=-it\,,
\]
and $C'_-$ is the small quarter-circle below the real axis, seen in the right
part of the figure. It is easy to check that in the limit $\epsilon\to0$ this
part of the integral yields $-F(0)/4$.

On the other hand the factor $(e^{2\pi i\tau}-1)^{-1}$ does not decay into the
upper half plane, but one can write
\[
\frac1{e^{2\pi i\tau}-1}=-1-\frac1{e^{-2\pi i\tau}-1}\,,
\]
where the second term does decay into the upper half-plane, so that
\[
I_+=-\int_{C_+}F(\tau)\dd\tau-\int_{\tilde C_+}\frac{F(\tau)\dd\tau}{e^{-2\pi i\tau}-1}
=\int_\epsilon^\infty F(\tau)\dd\tau+i\int_\epsilon^\infty\frac{F(it+0)\dd t}{e^{2\pi t}-1}+\int_{C'_+}\,,
\]
where the second term in the last form represents the contribution of the
straight part of $\tilde C_+$ (with $\tau=it+0$), and the last term is the
contribution of the small quarter-circle above the real axis; again, it is
easy to check that this last contribution is $-F(0)/4$ in the limit
$\epsilon\to0$. Adding $I_+$ to $I_-$, and taking the limit $\epsilon=0$,
yields \eqref{eq:AbelPlana}.

\begin{hwproblem}{2}\label{hw:2}
Derive Casimir force generated by massless scalar field theory in
$d+1$-dimensional space-time ($d$ spatial dimensions plus time). Assume that
the plates have $d-1$ spatial dimensions, i.e.\
$\xx=(\xx_\parallel,x_\perp)$, $\xx_\parallel\in\R^{d-1}$.
\end{hwproblem}
